\documentclass[12pt]{article}

% House style preamble
\input{.house-style/preamble.tex}
\usepackage{float}
\usepackage[section]{placeins}

% Update PDF metadata
\hypersetup{
  pdftitle={From Scores to Kinds: Projectibility and Validity in AI Evaluation},
  pdfkeywords={AI evaluation, benchmarks, construct validity, projectibility, inference licences, natural kinds}
}

\title{From Scores to Kinds: Projectibility and Validity in AI Evaluation}
\author{Brett Reynolds \orcidlink{0000-0003-0073-7195}%
\thanks{Contact: \href{mailto:brett.reynolds@humber.ca}{brett.reynolds@humber.ca}}\\
Humber Polytechnic \& University of Toronto}
\date{\today}

\fancyhf{}
\fancyhead[L]{\small\scshape From Scores to Kinds}
\fancyhead[R]{\small\thepage}
\fancyfoot{}
\renewcommand{\headrulewidth}{0pt}
\setlength{\headheight}{14pt}
\setlength{\emergencystretch}{2em}

\begin{document}
\maketitle

\begin{abstract}
AI evaluation terms let scores and labels travel beyond the cases that
produced them. A benchmark result becomes evidence for \term{reasoning}; an
attack result becomes evidence of \term{jailbreak resistance}; an aggregate
label becomes a governance claim about \term{alignment} or
\term{trustworthiness}. Such terms function as inference licences: they
authorize movement from an observation or score to a claim about recurrence,
vulnerability, capability, or governance. Evaluation practice has explicit
norms for two questions. Reliability asks whether a result is stable;
construct validity asks what it means in the measured setting. Neither says
when the licensed inference may travel, or what keeps it alive.

This paper supplies the structure of that third norm. Drawing on Goodman's
projectibility and on Boyd's and Khalidi's accounts of kinds, I treat each
particular evaluation as issuing a licence with a life-cycle. A licence is
generated under validity conditions, or blocked, or vitiated (contamination
is the clearest vitiator); it's bounded in scope; and once live it can be
narrowed by new evidence, eroded by optimization against the measure,
defeated by stronger results, revoked, or left to expire as models and
policies change. Two axes cut across every stage: misrecognition, when
uptake outruns what was licensed, and maintenance, since stability under
repeated measurement isn't control. A four-role typology of evaluation terms
falls out as a diagnostic, and worked cases (hallucination; reasoning and
emergence; jailbreak; alignment and trustworthiness, with sycophancy as
pressure case) show the life-cycle sorting disputes the labels alone can't:
the emergence debate, for one, becomes a definite question about when a
capability licence forms.
\end{abstract}

\medskip
\noindent\textbf{Keywords:} AI evaluation; benchmarks; construct validity;
projectibility; inference licences; natural kinds

\section{Introduction}\label{sec:introduction}

AI evaluation often moves from a score to a much larger claim. A benchmark result
is reported as evidence for \term{reasoning}. A safety evaluation says that a
system has \term{jailbreak resistance}. A model card, audit, or policy document
describes an assistant in terms of \term{alignment}, \term{trustworthiness}, or
\term{hallucination} risk. These labels point to observed behaviour and invite
readers to draw inferences: that a failure will recur, that a defence will
generalize, that a capability will transfer, or that a system is suitable for
some governance purpose.

The debate over emergent abilities in large language models shows the problem in
a compact form. \textcite{wei2022emergent} argue that some abilities appear only
once models reach sufficient scale. \textcite{schaeffer2023emergent} reply that
some apparent jumps can be produced by the choice of metric: discontinuous
scoring can make gradual changes in performance look like sudden emergence.

The dispute reaches beyond the acquisition of new abilities. It asks what a score
licenses us to infer. If a benchmark curve has a sharp bend, can we
infer a new capability, a change in the measurement instrument, or something
more local about a task, metric, and model family?

The problem is general: \textcite{raji2021whole}
argue that influential benchmarks are often treated as stand-ins for broad
claims about general ability, even though the benchmarks themselves are finite,
contextual, and constructed around specific tasks. A benchmark can be useful
while still being too narrow to support the claim attached to it. AI evaluation
terms are vague in places, but the deeper issue is inferential. They function as
\term{inference licences}: they authorize a move from an observation or score to
a claim about recurrence, vulnerability, capability, or governance. The idea has
a pedigree. Ryle treated law-like statements as inference tickets, statements
whose content is what they entitle us to infer rather than a hidden state they
describe \citep{ryle1949concept}. Different evaluation terms authorize
different moves.

This paper asks when such licences are earned, how far they travel, and what
keeps them alive. A term licenses a type of move; a particular evaluation
issues a particular licence; and the licences, not the terms, are what have
careers. Which move each term licenses is the diagnostic sub-question, and the
worked cases pursue it term by term.

I use \term{projectibility} in a plain sense first: a classification is
projectible when applying it in one case gives us a warranted reason to expect
something in another. The term is Goodman's, coined for the predicates that
sustain induction \citep{Goodman1955}; I extend it to the classifications an
evaluation practice deploys.

If a system produces one
unsupported answer, that alone doesn't show a general tendency to hallucinate.
If an attack succeeds once, that alone doesn't show that the model is generally
unsafe. If a model passes one reasoning benchmark, that alone doesn't show a
stable reasoning capability. The classification becomes useful when we can say
what else should follow from it, in which domain, and under which limits.

In a classical definition, a
term gives necessary and sufficient conditions for membership. Boyd's account
of \term{homeostatic property clusters} and Khalidi's causal-network account
show how classes can support projection without that form of definition
\citep{boyd1991,boyd1999,khalidi2018causal}. For AI evaluation, the lesson is
modest: useful terms needn't have essences, but they should mark structures that
support defensible inferences beyond the immediate observation.

The \term{construct validity} tradition gives the measurement side of the same
problem. A test score doesn't interpret itself. It needs a construct, an
operationalization, a metric, a population of instances, and an intended use.
Kane's argument-based view of validation is useful here because it treats
validity as a question about the interpretations and uses licensed by scores
\citep{kane2013validating}. AI evaluations should be read the same way. A label
such as \term{hallucination} or \term{jailbreak} belongs to a measurement
practice that includes tasks, prompts, raters or automated judges, policies,
metrics, populations, and deployment decisions.

The division of labour is the paper's hinge, and it has three parts rather than
the familiar two. Reliability asks whether a result is stable when the
measurement is repeated. Construct validity asks what a
score or label licenses in the measured setting. Projectibility asks when that
licensed inference can travel to another model, benchmark, attack, or
deployment, and what maintains it over time. Evaluation practice has explicit
norms and methods for the first two questions. For the third it mostly has
habits.

Neither the kind theory nor the validity theory is the contribution claimed
here. Boyd and Khalidi supply projection without essences; Kane supplies the
argument-based view of validation; Goodman supplies the vocabulary. The
contribution is a structure none of them supplies: a life-cycle for evaluation
licences. A licence is generated under stated conditions, or blocked, or
vitiated; it's bounded in scope; and once live it can be narrowed, eroded,
defeated, or revoked, or left to expire, while the community's uptake tracks
that career accurately or misreads it. The life-cycle turns projectibility from a
compliment paid to good benchmarks into a set of stages with stateable norms.

The argument proceeds as follows.
Section~\ref{sec:licences} presents evaluation terms as inference licences and
connects them to construct validity.
Section~\ref{sec:projectibility} introduces projectibility without essences and
distinguishes it from formal transportability.
Section~\ref{sec:lifecycle} sets out the licence life-cycle and restates the
paper's typology of evaluation terms as a diagnostic that falls out of it.
Sections~\ref{sec:hallucination} and~\ref{sec:capability} work through the two
central cases, hallucination and capability attribution; the second returns to
the emergence dispute and locates it at a definite stage of the life-cycle
rather than adjudicating it.
Section~\ref{sec:jailbreak} treats jailbreak more briefly, as the case where
licences age fastest. Section~\ref{sec:governance} recasts alignment and
trustworthiness as thin verdicts issued by coordination regimes, with
sycophancy as the pressure case. Section~\ref{sec:objections} answers
objections and concedes limits.

\section{Evaluation terms as inference licences}\label{sec:licences}

Evaluation language climbs an inference ladder. A raw output, interaction, or
score becomes an evaluation label. The label receives a construct
interpretation. The construct supports a kind attribution. The attribution then
enters a model comparison, release decision, procurement standard, or governance
claim. Each step can be legitimate, but the warrant changes as the claim climbs.

\begin{figure}[H]
\centering
\small
\begin{tabular}{@{}c@{}}
raw output, interaction, or score \\
\(\downarrow\) \\
evaluation label \\
\(\downarrow\) \\
construct interpretation \\
\(\downarrow\) \\
kind attribution \\
\(\downarrow\) \\
comparison, release, procurement, or governance claim
\end{tabular}
\caption{The inference ladder in AI evaluation.}
\label{fig:inference-ladder}
\end{figure}

Consider a mathematical-reasoning benchmark. Raw completions are first coded as
right or wrong on individual items. The metric aggregates those judgements into
an accuracy score. The benchmark report then interprets the score as evidence
for a construct such as \term{mathematical reasoning}. A leaderboard turns that
interpretation into a model ranking, and a user can then treat the ranking as
evidence for deployment suitability. The same number appears throughout, but the
claim has changed at every rung.

The narrow inference can be sound while the wider one fails. A high score on one
population of word problems can license a claim about performance on those
items, under that prompt format and scoring rule. It doesn't yet license a
general reasoning claim across paraphrases, irrelevant clauses, new item
generators, languages, tool scaffolds, or contamination controls. The ladder
makes the analytic question visible: which movement is being licensed,
and what evidence supports it?

Construct validity names the first warrant. A score has meaning only relative to
a construct, an operationalization, a metric, a population of instances, and an
intended use \citep{kane2013validating,chouldechova2024shared}. An
operationalization is the concrete test procedure that makes a construct
measurable. Here, validity is local: the relevant warrant connects a score to a
stated interpretation and use. Work on LLM benchmarks makes the same point in
the evaluation setting: a benchmark operationalizes a phenomenon through tasks
and metrics, then supports a claim only if those parts line up
\citep{bean2025measuring,freiesleben2026establishing}.

In LLM evaluation, repeated probing often makes \term{reliability} easier to
estimate: evaluators can ask whether results are stable across item samples,
prompt variants, decoding settings, and evaluator systems. Reliability sharpens
the local warrant, but it doesn't settle interpretation. The validity question is
which probes are compressed into which score, for which construct interpretation
and use.

Model comparison makes that compression consequential. Ranking model A above
model B requires more than running the same benchmark twice. It requires a
judgement that the chosen score summarizes the right profile of tasks, prompts,
languages, evaluator systems, and deployment conditions for the intended
comparison. A leaderboard can be valid as a narrow benchmark comparison and fail
as a governance conclusion if those supports don't hold \citep{raji2021whole}.

A licence isn't the same as its uptake. When a leaderboard, a model card, or a
procurement office treats a score as supporting a claim, that's uptake: real,
consequential, and sometimes wrong. What the score licenses is fixed by its
generation conditions and scope, not by how confidently it's repeated.
Mistaken uptake has effects (deployments, purchases, headlines), but it doesn't
widen the licence, any more than an official record of a permission widens the
permission. The gap between licence and uptake isn't noise. It's one of the
framework's objects of study, and Section~\ref{sec:lifecycle} gives it a name.

Model comparison also marks the limit of construct validity. Validity can
say what the score licenses here, but it doesn't settle whether the claim
travels to another model family, attack distribution, user population, or
deployment. That's the projectibility question.

\section{Projectibility without essences}\label{sec:projectibility}

The projectibility question turns on structure rather than classical
definitions. A classical definition gives necessary and sufficient conditions
for class membership. The accounts I draw on reject that requirement. Boyd's
account of \term{homeostatic property clusters} explains how a class can
support projection when its properties cluster through stabilizing mechanisms.
Khalidi's causal-network account explains projectibility through relations among
properties, where some properties are more central and help explain the presence
or stability of others \citep{boyd1991,boyd1999,khalidi2018causal}.

Here, \term{class} and \term{kind} do different jobs. A class is any grouping
introduced by an evaluation practice. A kind, in the stronger sense used here,
is a class whose structure supports projection. The paper doesn't ask
whether every label is a kind. It asks when an evaluation class earns that
stronger status. The stipulation isn't circular, because the order of analysis
runs from structure to projection: the worked cases identify the structure
first and derive the licensed inferences from it.

A classification is projectible when applying it in one setting gives a
warranted reason to expect something in another setting, under stated
conditions. The conditions matter. Projection is field-relative: what travels
depends on the domain, purpose, population, and use of the evaluation. A label
can project across nearby prompt variants while failing to project across
languages, tool access, deployment settings, or adversary classes.

Suppose a red-team test finds that an adversarial suffix (a string appended to a
prompt) makes a chat assistant answer a prohibited request. The result supports
a narrow projection. If the suffix still works after a few tokens are swapped,
after a harmless phrase is added before it, or when another optimizer seed
produces a similar appended string, the test gives reason to expect similar
vulnerability under the same chat interface, refusal policy, and evaluator
standard.

The same result says little about an indirect prompt-injection case where a
browser agent reads a web page that tells it to send a private file. The
attacker, channel, policy boundary, and system boundary have changed, so the
original classification no longer carries the same expectation.

In AI evaluation, projectibility-relevant structure is artificial, engineered,
and institutional. It can include training data, model architecture, training
objectives, post-training incentives, decoding procedures, retrieval design,
prompt format, interface constraints, policy boundaries, evaluator models,
benchmark construction, reporting conventions, and deployment use. None of these
is necessary or sufficient; none is essential. Together, though, they can explain
why a failure profile, vulnerability, capability attribution, or governance
bundle recurs, and why projection weakens when those supports change.

The travel question already has a formal treatment. Bareinboim and Pearl's
transportability
results give graphical conditions under which causal and statistical relations
estimated in one setting carry to another \citep{bareinboim2016causal}. The
formalism presupposes two things the evaluation case usually lacks. The
variables are fixed in advance, and the differences between source and target
are encoded in a causal model over those variables. Projectibility, as used
here, is the prior question: which classifications deserve to be treated as
variables at all. That's Goodman's question, not a transport question, and no
one can write a selection diagram over incentive structure, refusal machinery,
and evaluator standard. Transport is also silent about use: nothing in it
distinguishes the same score feeding a release gate from the same score feeding
a procurement rule. Both gaps matter for evaluation, and both are where the
kind-theoretic accounts do their work.

A structure matters for projectibility when it does two jobs: it explains why an
inference travels and where it stops. A resampling detector gives the simple
case. The detector treats variation across sampled answers as evidence that a
claim is unsupported. The signal can travel to settings where unsupported claims
vary across samples. It travels poorly to a model that repeats the same
unsupported answer every time, because the signal has disappeared.

Structure of this sort explains travel at a time. But licences also live in
time: they're issued, worn down, overtaken, and withdrawn. The next section
sets out that life-cycle, and the worked cases then apply both questions
together: what structure lets an inference travel, and what stage of its career
the licence has reached.

\section{The life-cycle of an evaluation licence}\label{sec:lifecycle}

An inference licence isn't issued once and for all. It's generated under
conditions, bounded in scope, and then kept, narrowed, worn out, overtaken, or
withdrawn. The shape is familiar from everyday normative statuses such as
permissions, which are granted under conditions, bounded, and later revoked,
exhausted, or exceeded. Nothing metaphysical rides on the analogy. Every stage
below names something evaluation practice already does; what the life-cycle
adds is order, assigning each failure its own name instead of the one word
\mention{invalid}.

\subsection{Generation and scope}

A licence is generated when the validity argument goes through: a construct, an
operationalization, a metric, an item population, and an intended use that line
up \citep{kane2013validating,bean2025measuring}. Generation has three outcomes,
not two.

In the central case the licence forms. The evaluation supports its stated
interpretation, and the licence's scope is fixed by what was actually sampled
and held constant: the item population, prompt format, evaluator standard,
system boundary, and time of measurement.

Sometimes nothing forms. Benchmark contamination is the clearest case: if the
test items sat in the training data, a perfect score licenses nothing about the
capability, because what was measured is closer to retrieval of the answer key.
The score is real; the licence is blocked. Blockage is an absence, not a weak
licence.

Between these lies \term{vitiated generation}: something forms, but less than
the report claims. Partial contamination, an unvalidated automated judge, and
unmeasured prompt sensitivity each leave a licence standing that's narrower or
weaker than its advertised scope. Many disputes about whether a benchmark
\enquote{really measures} a construct are disputes about vitiation, and they go
better when put that way, because vitiation turns on specifiable conditions
rather than on essence.

A licence used beyond its scope isn't a weaker
licence; for the wider claim, it's no licence at all. A source-faithfulness
score in retrieval-augmented generation licenses claims about that corpus and
retrieval design, not about open-domain factual accuracy. Exceeding scope is
the characteristic misstep that the misrecognition axis below will name.

\subsection{The career of a live licence}

Five things can become of a licence once generated.

New evidence can \term{narrow} it. When controlled variants show that
performance drops under value substitutions and irrelevant clauses, the
original licence isn't cancelled; it persists with reduced scope, covering the
original items and format but no longer the capability reading
\citep{mirzadeh2025gsm}. A narrowed licence re-enters the cycle, where it can
be narrowed again, defended, or given up.

Optimization against the measure \term{erodes} it. A benchmark licence is
issued under a tacit no-targeting condition: the score means what it means
partly because no one was steering toward it. Reward-model overoptimization
gives the general form of the problem, since optimizing a learned proxy outruns
the target it stood in for \citep{gao2023scaling}, and leaderboard culture
applies the same pressure to benchmarks. The familiar slogan that a measure
which becomes a target stops being a good measure has a licence-theoretic
reading: targeting doesn't falsify the score, it spends the licence. Saturation
is erosion's endpoint.

The issuing conditions can lapse, and the licence \term{expires}. Models are
updated, data distributions move, policies change. An evaluation of last
year's model licenses nothing about this year's release, but the number
survives in citation graphs and marketing long after the configuration it
measured has gone. Expiry is undramatic, which is why it's the easiest stage to
miss.

A stronger result can \term{defeat} it without invalidating it. A red-team
success can defeat a deployment licence while the benchmark result it defeated
still stands as an accurate record of what it measured. Defeat is override, not
refutation, and the difference matters: a defeated licence can return if the
defeating result is itself narrowed or expires.

The issuing practice can \term{revoke} it. Benchmarks are retracted when
contamination is discovered, automated judges are replaced when bias is found,
datasets are withdrawn. Revocation is uptake-facing: it tells the community to
stop inferring, which is different from a licence quietly expiring.

\subsection{Misrecognition and maintenance}

One axis cuts across every stage. Each transition can be tracked, missed, or
mistaken by the community that uses the licence. Uptake can treat a vitiated
licence as full, an expired licence as live, a narrowed licence at its original
scope. Model cards and surveys repeat claims whose supports have lapsed, and
repetition doesn't reissue a licence. \term{Misrecognition} is second-order: a
fact about the relation between a licence and its uptake, and explanatory in
its own right. Benchmark essentialism, the habit of reading a benchmark as the
capability itself, isn't a one-off error but an entrenched misrecognition
pattern with institutional supports: leaderboards reward it, citation practices
spread it, and procurement shorthand depends on it.

The last piece is \term{maintenance}, where stability and control come apart.
Repeated measurement can show that a score is stable; it can't keep a licence
current, any more than re-reading a thermometer maintains the temperature.
Weinberger's distinction between homeostasis and mere stability makes the point
general: a system is controlled only when something corrects perturbations
\citep{weinberger2026HomeostasisCausalControl}. An evaluation ecosystem
maintains licences only where it has that control structure: monitoring that
notices drift, re-evaluation under new items and attacks, release gates that
respond to defeat, reissue after expiry. A leaderboard that never retests is a
record, not a controller, and the licences it displays age at the speed of the
field.

\subsection{A diagnostic typology}\label{sec:typology}

Single claims have careers. Across many claims, recurring term-roles show
characteristic careers, and a compact typology is useful as a diagnostic.
Table~\ref{tab:typology} classifies AI evaluation terms by the inference their
evaluations are mainly used to license. The classification rule is a
stipulation about primary role, not a discovery about usage statistics, and
nothing below depends on a borderline assignment. Read each row from left to
right: an observation becomes an operationalization, the operationalization
licenses an inference, and the final columns state what can travel and which
misrecognition the term characteristically invites.

\begin{table}[H]
\centering
\scriptsize
\setlength{\tabcolsep}{2pt}
\renewcommand{\arraystretch}{1.15}
\begin{tabular}{>{\raggedright\arraybackslash}p{0.11\textwidth}
                >{\raggedright\arraybackslash}p{0.09\textwidth}
                >{\raggedright\arraybackslash}p{0.14\textwidth}
                >{\raggedright\arraybackslash}p{0.15\textwidth}
                >{\raggedright\arraybackslash}p{0.15\textwidth}
                >{\raggedright\arraybackslash}p{0.13\textwidth}
                >{\raggedright\arraybackslash}p{0.10\textwidth}}
\toprule
Type & Observation & Operationalization & Licensed inference & Projectibility payoff & Characteristic misrecognition & Use context \\
\midrule
Output-centred failure profile &
generated output &
factuality, faithfulness, support, uncertainty, granularity &
the failure is likely to recur under specified structural supports &
predicts recurrence, detection, and mitigation transfer under similar generation
or grounding conditions &
treating \term{hallucination} as false output, or treating one low score as
deployment safety &
model evaluation, retrieval-augmented generation, factuality claims \\
\addlinespace
Adversarial relation &
boundary-crossing interaction &
attack distribution, policy class, evaluator, target system &
the system is vulnerable under a specified adversarial network &
predicts vulnerability and defence brittleness relative to attacker class,
interface, and policy boundary &
treating \term{jailbreak resistance} as a model property or projecting across
attacker classes &
red-teaming, release gates, security evaluation \\
\addlinespace
Capability-attribution term &
task performance &
benchmark, metric, item population, prompting scaffold &
the performance supports a capability claim over a domain &
predicts transfer or robustness across item variants, metrics, scaffolds, and
system boundaries when the attribution is validated &
treating a score as a capability, or a benchmark as the domain &
benchmark reporting, model comparison, deployment claims \\
\addlinespace
Umbrella governance construct &
aggregate risk or value label &
component taxonomy, governance use &
lower-level constructs are being coordinated for a purpose &
predicts which lower-level evaluations should travel together for a stated
governance decision &
treating \term{alignment} or \term{trustworthiness} as a single behavioural
construct with a single score &
procurement, audit, release and risk-management decisions \\
\bottomrule
\end{tabular}
\caption{AI evaluation terms by primary inferential role.}
\label{tab:typology}
\end{table}
\FloatBarrier

The typology is non-exclusive. Hallucination has relational features because
unsupportedness depends on a source, a task, and a user context; jailbreaks
often end in bad outputs; agency can be described as a capability, a system
architecture, or a deployment relation. A term's row records its primary role,
nothing more.

Confused debates arise when one row's licence is read with another row's
scope: a score read as a capability, a relational vulnerability read as a model
property, an umbrella verdict read as a single behavioural cluster. The worked
cases make the diagnostic concrete, and each also displays a different part of
the life-cycle. Hallucination shows structure and scope; capability
attribution shows generation, narrowing, and erosion; jailbreak shows expiry
and defeat; and
the governance verdicts show what uptake looks like when regimes rather than
researchers hold the licences.

\section{Hallucination: the kind behind the licence}\label{sec:hallucination}

\term{Hallucination} starts from a generated output, but the failure profile is
wider than falsity. It includes factuality, faithfulness to source, support and
grounding, uncertainty, granularity, retrieval conditions, task context, and user
uptake. The output remains the centre of the evaluation, but the source,
retrieval system, prompt, and use context help determine which failure profile the
output instantiates \citep{ji2023survey,huang2025survey}.

The source relation is central. In source-conditioned generation, an output can
contradict the source, add claims that the source doesn't verify, or introduce
content that's true but unsupported by the assigned evidence. In open-domain
question answering, the relevant source can be world knowledge rather than an
input document. In retrieval-augmented generation, it can be the retrieved
context, the retrieval corpus, or a cited passage. These settings share a family
resemblance, but they don't license the same inference.

False output, then, is too coarse. A model can make a factual error,
overstate what the available source warrants, contradict a retrieved document,
fail to retrieve the needed evidence, answer with misplaced confidence under
uncertainty, or give an answer whose harm depends on the user's task. These are
related failures, but their operationalizations differ. A factuality metric for
biographical claims doesn't measure source faithfulness in summarization, and a
source-faithfulness score in retrieval-augmented generation doesn't measure
whether a medical or legal assistant should have abstained.

Task goal matters too. If a system is asked to write a true-or-false test and
includes \enquote{Toronto is the capital of Canada} as an item, the false
proposition is intentional test content, not a hallucination. The same sentence
becomes a hallucination if the system presents it as true, treats it as
supported, or uses it as an answer where accuracy is the goal. The label tracks
the role the output is meant to play, not falsity in the string alone.

The failure profile has
recurring parts: a plausible generated answer, an evidential or factual defect,
some relation to a source or task goal, and a user-facing risk created by the
answer's uptake. Different evaluations put weight on different parts of the
profile. A biography generator invites atomic factual checks; a summarizer
invites source-faithfulness checks; a high-stakes assistant invites abstention
and uncertainty checks.

So far, that's a list, and a list can't say which projections the
classification licenses. The parts of the profile are ordered: some generate
others, and the ordering is where the licence gets its content
\citep{ji2023survey,huang2025survey}. Gaps in parametric knowledge generate the
factual-error profile. Decoding and sampling choices generate the fluency and
stated confidence that mask unsupportedness and give the failure its
characteristic plausibility. Training and evaluation regimes that reward
answering over abstaining generate the calibration profile, converting
uncertainty into confident assertion. Retrieval design generates the
source-faithfulness profile, because it fixes which evidence the output can be
unfaithful to \citep{shuster2021retrieval}. In Khalidi's terms, the profile is
a node with internal causal structure, not a checklist
\citep{khalidi2018causal}.

That ordering explains the shape of the measurement literature. FActScore
decomposes long-form answers into atomic facts and asks whether a specified
source supports them: it reads the source relation \citep{min2023factscore}.
SelfCheckGPT treats instability across sampled responses as a black-box signal
of unsupported content: it reads the decoding position
\citep{manakul2023selfcheckgpt}. Retrieval augmentation intervenes on the
evidence boundary rather than reading it, and can introduce noisy, missing, or
conflicting context of its own \citep{shuster2021retrieval,huang2025survey}.
Uncertainty-based methods read the calibration position: under some
conditional-generation settings, higher predictive uncertainty tracks greater
hallucination risk \citep{xiao2021hallucination}.

Each operationalization licenses projection along the part of the structure it
reads, and fails where its position predicts. A resampling detector,
Section~\ref{sec:projectibility}'s simple case, misses the model that repeats
the same false claim in every sample, because the decoding signal has
disappeared. An atomic-fact score misses relevance, recall, and
whether the system should have answered at all, because those sit away from
the source relation. A retrieval intervention helps where missing parametric
knowledge was the bottleneck and adds failure modes of its own where it wasn't.
These boundaries aren't peripheral complications. They're the licence's scope
conditions, read off the kind's structure.

Read as a licence, then, a hallucination classification supports projection to
recurrence and mitigation transfer where the generating structure is shared:
the same knowledge boundary, decoding procedure, retrieval design, and
incentive profile. Change one generator and the projection weakens exactly
there, which is what a structured kind predicts and a flat label doesn't. And
no position in the structure licenses the jump from one low hallucination score
to deployment safety. That inference needs the coordination machinery of
Section~\ref{sec:governance}, stated rather than assumed.

\section{Capability licences: reasoning, emergence, and agency}\label{sec:capability}

\term{Reasoning}, \term{emergence}, and \term{agency} share a score-to-capability
problem. Each begins with task performance and asks whether that performance
supports a claim about what the system can do across a domain. The failure
routes differ. Reasoning claims can be fragile across item variants, distractors,
scaffolds, and contamination controls. Emergence claims can depend on metric
choice. Agency claims can depend on where the model-tool-environment boundary is
drawn.

The common risk is treating the benchmark as if it directly revealed the
capability. A task score is an observation, a benchmark aggregate is an
operationalization, and a capability attribution is an interpretation. The
attribution becomes projectible only when the benchmark is embedded in a wider
network of evidence: item variation, tests that separate the target construct
from nearby constructs, contamination checks, prompting conditions, and
comparisons with neighbouring constructs \citep{freiesleben2026establishing}.

In life-cycle terms, the recurring troubles sort by stage. Contamination is a
generation failure: the capability licence never forms, however high the score.
Unmeasured prompt sensitivity vitiates generation. Scaffolds set scope: if
chain-of-thought prompting supplies the deliberation, the licence covers
scaffolded performance, not unaided problem solving. And controlled item
variation narrows licences after issue. Contamination checks are generation
conditions, not optional extras.

Reasoning is the central case. The survey literature treats reasoning as
a plural and unsettled construct, covering deductive, inductive, abductive,
analogical, causal, probabilistic, formal, and informal reasoning
\citep{huang2023reasoning}. That plurality matters for validity. A benchmark can
show strong performance on one task population while leaving open whether the
model has a reasoning capacity that transfers across problem forms, domains, and
prompting regimes.

The plurality also changes what counts as confirming evidence. A deductive
reasoning claim asks for validity-preserving steps across form changes. A causal
reasoning claim asks for sensitivity to interventions and counterfactual
structure. A mathematical word-problem claim asks for stable mapping from the
story to the relevant operations. Treating these as one flat reasoning score
makes the construct easier to report and harder to validate.

GSM-Symbolic shows narrowing in action. \textcite{mirzadeh2025gsm} generate
controlled variants of GSM8K-style problems and show that model performance can
vary substantially across value substitutions, problem complexity, and seemingly
irrelevant clauses. The original licences aren't cancelled; they're narrowed,
still covering the original items under the original format but no longer the
general capability reading the reports had given them. A reasoning attribution
grows stronger the more narrowing it survives: controlled variation,
contamination checks, distractors, and changes of scaffold.

Emergence shows the apparatus situating a dispute rather than adjudicating it.
\textcite{wei2022emergent} issue capability licences on an operational
condition: an ability present in larger models and absent in smaller ones.
\textcite{schaeffer2023emergent} argue that under nonlinear or discontinuous
metrics the generation conditions were never met: a sharp bend in the curve can
report the metric's geometry rather than a new capability. Put this way, the
disagreement stops being a clash over the word \mention{emergence} and becomes
a definite question about generation: which invariances, across metrics, item
populations, and model families, have to hold before a capability licence
forms, and which are merely desirable. The apparatus doesn't pick a side. It locates
the fault-line, and a framework that can situate a dispute does more than one
that can only label its parties.

The constructive lesson is to report the metric and task structure that make an
emergence claim projectible. A sharp threshold under one scoring rule can license
a claim about performance under that rule. It licenses a broader capability
claim only when neighbouring metrics, item populations, and model families
support the same interpretation.

Benchmark infrastructure then reads as licence maintenance. MMLU's breadth
widens scope, but its lopsided per-subject performance and calibration problems
warn against reading the aggregate at face value \citep{hendrycks2020measuring}.
HELM widens the generation conditions, requiring calibration, robustness,
fairness, bias, toxicity, and efficiency alongside accuracy before broad
licences issue \citep{liang2023holistic}. Dynabench is reissue machinery: when
static benchmarks saturate, which is erosion's endpoint,
human-and-model-in-the-loop collection generates new items and, with them, new
licences \citep{kiela2021dynabench}.

Agency is the system-boundary case, where everything turns on scope. Agentic LLM
surveys define the target as a system that receives natural-language input,
reasons or plans, acts autonomously, uses tools or environments, and pursues
goals \citep{plaat2025agentic}. The attribution attaches to a
model-system configuration, not to a bare language model. Tool access,
memory, planning horizon, feedback, permissions, and environment design help
determine whether the system counts as agentic for a given evaluation.

The same base model can be part of different agency profiles. In an autocomplete
interface, it predicts text without acting in the world. In a tool-using
assistant, it can call APIs, revise plans, and affect an external state. In a
multi-agent simulation, its behaviour depends on other agents and feedback from
the environment. The agency licence travels only with the relevant system
boundary.

\term{AGI} is the limiting case of a capability attribution. I've argued
elsewhere that AGI evaluation goes better when general intelligence is treated
as a property cluster rather than a checklist of tasks \citep{reynolds2025agi}.
In the present terms, an AGI attribution is a capability licence whose scope is
contested all the way down, which is why no single benchmark could generate it.

The question \enquote{Does the model reason?} then dissolves into stages. What
licence was generated, over which item population and scaffold? How far has it
since been narrowed? How much erosion has leaderboard pressure inflicted? A
capability term earns its keep when the answers predict transfer, robustness,
and failure boundaries better than the raw score alone.

\section{Jailbreak: licences under an adaptive adversary}\label{sec:jailbreak}

\term{Jailbreak} starts from an interaction that crosses a policy boundary. The
evaluation object is a relation among model behaviour, attacker strategy,
interface affordances, guardrail design, evaluator judgement, and the
institutional definition of prohibited output. That's why jailbreak resistance
should be paired with an adversary model rather than treated as an intrinsic
property of a model \citep{yi2024jailbreak,vassilev2024adversarial}. Direct
chat jailbreaks, transferable adversarial suffixes, persuasion-based attacks,
and indirect prompt injection share enough of this relational profile to be
grouped together, but they don't share one mechanism
\citep{zou2023universal,zeng2024johnny,greshake2023not}.

This relational structure explains why in-the-wild jailbreaks are a social and
technical population. Prompt families are copied, revised, hosted, tested, and
patched across time. Their success depends on vendor policy, model update,
interface, forbidden-content class, and evaluator choice
\citep{shen2024do}. A benchmark that reports attack success rate without this
network can still be useful, but the reported number has a narrow inferential
target.

What distinguishes this row of the typology is tempo. Capability licences erode
slowly; jailbreak licences date fast, because one of their scope conditions,
the attack distribution, is an adversary that adapts on purpose. A resistance
claim can be defeated by the next optimizer run and expires with the next
interface or policy change. That's why red-teaming can't be a pre-release
ritual. In life-cycle terms it's maintenance: the control loop that reissues or
revokes resistance licences as the attack population moves
\citep{weinberger2026HomeostasisCausalControl}. A safety case that cites last
quarter's attack suite is running on an expired licence.

Classifying jailbreak as an adversarial relation licenses projection only
relative to a specified adversarial network. A result can support expectations
about vulnerability under attacker class A, interface I, policy boundary P, and
attack distribution D. It doesn't by itself support expectations about attacker
class B, a different interface, or a new policy regime. The payoff is a more
precise release claim: resistance is always resistance to a defined attack
space, for as long as that space stays the relevant one.

\section{Alignment and trustworthiness: thin verdicts and coordination regimes}\label{sec:governance}

\term{Alignment} and \term{trustworthiness} differ from the other rows in
role, not just in breadth. The typology called them umbrella constructs,
naming the bundle of lower-level evaluations they manage. Seen from the
label's side, they're \term{thin verdicts}: they mark the outcome of an
assessment without giving the route through the evaluations that produced it,
the way \mention{wrong} marks a moral conclusion without saying what was wrong
with the act. A thin verdict projects only after some regime has fixed which
lower-level evaluations feed it and how. Treating these terms as first-order
behavioural kinds misfires: the projectible structure sits in the
regime, not in the behaviour.

NIST's AI Risk Management Framework is best read in exactly that way, as a
constitution for a coordination regime. It bundles validity and reliability,
safety, security and resilience, accountability and transparency,
explainability and interpretability, privacy, and fairness with harmful bias
managed into \term{trustworthy AI}, for the stated purposes of risk management
\citep{nist2023artificial}. The bundle is the content.

A \term{coordination regime}, as I'll use the term, is the institutional
structure that takes evaluations up: its assignment rules (which results count,
at what thresholds), its records, its sanctions, and its decision points. A
release gate, a procurement rule, an audit, and a research leaderboard are
different regimes. The same lower-level result projects differently under each,
because each fixes a different bundle and a different decision. A release
decision and a procurement decision can use the same umbrella word while
licensing different projections. That's regime-relativity, not equivocation,
and it becomes equivocation only when the regime is left unstated.

The projectible claim is coordinative rather than behavioural. A
\term{trustworthiness} label predicts that a specified set of lower-level
evaluations has been assembled for a decision, that the same bundle can be
inspected again in a similar use, and that omitting one component changes the
governance claim. In life-cycle terms, the umbrella licence is only as current
as the stalest licence in its bundle: a trustworthiness verdict resting on an
expired robustness evaluation has expired too, whether or not anyone has
noticed.

\term{Alignment} is similarly target-dependent. One target is user preference:
RLHF trains models to follow instructions by optimizing against preference data
and reward models \citep{ouyang2022training}. Another target is truthfulness. A
third is policy compliance. A fourth is longer-term user welfare or institutional
risk management. These targets can reinforce one another, but the construct
claim changes when the target changes.

Consider a deployed assistant that gives agreeable answers to user beliefs,
corrects false factual claims, and refuses prohibited requests under a content
policy. A preference evaluation, a truthfulness evaluation, and a
policy-compliance evaluation can assign different scores to the same system.
Calling the system aligned says little until the lower-level target is stated.

Sycophancy is the pressure case. A behaviour can satisfy one alignment target,
such as immediate user preference, while failing another, such as truthfulness.
The same overoptimization that erodes benchmark licences
(Section~\ref{sec:lifecycle}) operates here inside the reward signal:
optimizing a learned proxy outruns the target it stood in for
\citep{gao2023scaling}. Sycophancy research makes the conflict concrete:
preference data and preference models can reward responses that match user
beliefs over truthful corrections \citep{sharma2024towards}, and recent formal
work argues that preference-based post-training can amplify this behaviour when
the learned reward favours agreement \citep{shapira2026rlhf}.

The pressure case matters because it can reverse the evaluative verdict without
changing the surface behaviour. Agreement with a false user belief can be a
success relative to a learned preference signal and a failure relative to
truthfulness, calibration, or epistemic support. The choice of lower-level
target determines which inference is licensed, and the thin verdict shouldn't
hide that choice.

The question \enquote{Is sycophancy misalignment?} is indexed to the
lower-level construct in view. Relative to a narrow preference target, some
sycophantic behaviour can be a local alignment success. Relative to truthfulness,
calibration, user autonomy, or epistemic support, it's a failure. Only a regime
that states its lower-level targets issues a determinate licence; an alignment
claim should say which target is in view and block projection to incompatible
targets. Without that specification, the thin verdict makes a decision sound
measured while leaving the measured construct underdescribed.

\section{Objections and limits}\label{sec:objections}

Three objections deserve answers, and two limits deserve concession.

The first objection says the licence vocabulary is metaphor. Licences are
deontic items, issued and revoked by authorities; evaluation warrants are
epistemic, and nobody signs them. The framework borrows the vocabulary for its
structure, not its metaphysics. Every stage names an event evaluation practice
already exhibits without the vocabulary: benchmarks are retracted, results are
patched around, claims outlive the configurations they measured, model cards
repeat lapsed numbers. What the deontic idiom adds is articulation: separate
names for failures that \mention{invalid} runs together, and a distinction
(licence versus uptake) that the idiom of truth and evidence doesn't draw by
itself. A reader who wants the idiom cashed out can read \mention{licence} as
\mention{warranted inference under stated conditions} throughout; nothing in
the worked cases changes.

The second objection says the framework proves too much: any evaluation dispute
can be redescribed as some stage of some licence's career, so the apparatus
explains everything and forbids nothing. But the account has failure
conditions. It would fail if evaluation claims that travel well systematically
lacked identifiable generation conditions and structural supports, since then
projectibility wouldn't track structure. It would fail if stage assignments
could be reshuffled case by case to save a diagnosis, since then the stages
would carry no information. And it would fail if licence and uptake never came
apart, since then the framework's central distinction would idle. The worked
cases put content into each condition: contamination is a generation fact,
GSM-Symbolic narrowing is an evidential fact, and benchmark essentialism is an
observed gap between what scores license and how they're used.

The third objection says the framework is external validity renamed. The reply
compresses Section~\ref{sec:projectibility}. External validity and
transportability ask whether a claim holds elsewhere, taking the variables and
categories as given \citep{bareinboim2016causal}. The present question is
prior, asking which classifications deserve projection at all, and wider,
asking what generates, maintains, and revokes the entitlement to project, and
who misreads it. A theory of truth-travel has no stage called revocation and no
axis called misrecognition; evaluation practice visibly has both.

Two limits should be conceded. The paper works two cases in depth and places
the others: jailbreak and the governance verdicts get sketches whose details
would need their own papers. And the framework classifies rather than
adjudicates. Locating the emergence dispute at generation doesn't settle it,
and mapping a licence regime doesn't certify the decisions the regime serves; a
well-maintained licence architecture can still serve a bad policy. Explanation
and vindication are different achievements, and the framework offers the first.

What remains is the contribution. Goodman supplies the term, Boyd and Khalidi
the structure-first account of projection, Kane the argument-based view of
validation, Weinberger the difference between stability and control. None of
them supplies the life-cycle: generation, scope, narrowing, erosion, expiry,
defeat, revocation, misrecognition, maintenance. The materials are borrowed;
the life-cycle is not.

\section{Conclusion}\label{sec:conclusion}

AI evaluation has explicit norms for reliability and validity, and mostly
habits for travel. The proposal here is that the missing norm has a structure.
Evaluation terms function as inference licences, particular evaluations issue
them, and the issued licences have life-cycles: each stage carries stateable
conditions, with misrecognition and maintenance as the cross-cutting axes. The
typology of term-roles then works as a diagnostic: output-centred failure
profiles, adversarial relations, capability attributions, and thin governance
verdicts fail in characteristically different ways, and disputes go wrong when
one role's licence is read with another's scope.

The worked cases show what the apparatus buys. Hallucination detectors turn
out to read different positions in one generating structure, which fixes where
each detector's licence stops. The emergence dispute relocates to a definite
question about generation conditions rather than a quarrel over a word.
Jailbreak resistance dates as fast as its adversary adapts, which makes
red-teaming maintenance rather than ritual. And an umbrella verdict is only as
current as the stalest licence in its bundle.

The framework has a practical corollary. An evaluation report could carry a
licence statement: the claim licensed, its scope, the structural supports it
rests on, the known breakers, and the conditions under which it should be
treated as expired. Model and system cards already have slots this information
would fit. Most of it exists at evaluation time and is lost by reporting
conventions built around single numbers; a licence statement is a format for
keeping it.

The typology is a starting point rather than a closed map. \term{Bias} is the
obvious extension. It looks partly like an output profile, partly like a
deployment-sensitive relation among model, population, task, and affected
group, and partly like a governance construct, so its licences will have an
unusually layered career. The useful question stays the same throughout: what
does the term license in a given evaluation, what structure lets the inference
travel, what would spend or defeat it, and who's checking that it's still
alive?

\section*{Acknowledgements}

The large language models OpenAI Codex, GPT-5.5 Pro, and Claude Fable 5 served
as drafting and editing aids throughout the preparation of this paper. I am
responsible for all theoretical claims, arguments, errors, and interpretive
choices.

\clearpage
\printbibliography

\end{document}
